\documentclass[10pt,a4paper]{amsart}
\usepackage[margin=19mm]{geometry}
\usepackage{amsmath,amssymb,mathtools}
\usepackage[scr=rsfs,cal=euler]{mathalfa}
\usepackage{xfrac}
\usepackage[unicode,hidelinks,bookmarks=false]{hyperref}
\newcommand{\ie}{i.e.\ }
\newcommand{\cf}{cf.\ }
\newcommand{\resp}{respectively\ }
\newcommand{\Subsection}{\subsection}
\providecommand{\nobreakdash}{\nobreak\hbox{-}}
\setlength{\emergencystretch}{2em}
\multlinegap0pt
\usepackage{bm}
\usepackage{bbm}
\usepackage{dsfont}
\usepackage{xspace}
\usepackage{cancel}
\usepackage{mathtools}
\usepackage{amsthm}
\makeatletter
\@ifundefined{lemma}{\newtheorem{lemma}{Lemma}}{}
\makeatother
\title[Assortment Optimization for Patient-Provider Matching]{Assortment Optimization for Patient-Provider Matching}
\author{Naveen Raman, Holly Wiberg}
\date{Source: arXiv:2502.10353v2 (CC BY 4.0)}
\begin{document}
\maketitle

\noindent\emph{Source and licence.} Assortment Optimization for Patient-Provider Matching. Version arXiv:2502.10353v2 (CC BY 4.0), distributed under the Creative Commons Attribution 4.0 International licence, as recorded in the arXiv licence metadata for this exact version and shown on the abstract page https://arxiv.org/abs/2502.10353v2. Full rights notice: licenses/arxiv-2502.10353-CC-BY-4.0.txt.\par\smallskip

\noindent\textit{Editorial scope.} Counted unit: the Lemma whose source label is \texttt{lem:augment} in the article (source0.tex lines 1126--1136, source label \texttt{lem:augment}), delivered together with the article's own complete original proof of it (source0.tex lines 1137--1166, one \texttt{proof} environment of 30 lines). No mathematics is written, repaired or re-derived: statement and proof are the article's own text line for line. Required context retained uncounted: none. Every definition, display and label of those lines is retained unchanged. Omitted from this package and not redistributed: 1--1125, 1167--1467. Nothing inside a retained excerpt is omitted or altered: every retained body is delivered exactly as the article's lines are written, so no omission receipt of any kind accompanies this package. Citations: the retained excerpts carry no citation command at all, so no bibliography entry of the article is transcribed and no reference is redistributed. Reference adaptation: none was needed and none was made; every reference inside the delivered unit resolves inside it, so no unresolved reference or citation is produced. Printed slips kept literal: no printed word, symbol, hypothesis or number of the counted statement or of its proof was altered. Layout and class: the article's own class is \texttt{article} with packages including geometry, authblk, hyperref, url, booktabs, amsfonts, nicefrac, makecell, microtype, xcolor, amsmath, amsthm, thmtools, thm-restate; the wrapper is the lane's pinned \texttt{amsart} editorial wrapper at 10pt on A4, which restates the theorem-like environments the retained text needs and adds the article's own macro definitions that the retained text needs (none), with extra wrapper packages (bm, bbm, dsfont, xspace, cancel, mathtools). The delivered numbering is the unit's own. Assets: no retained span contains a figure environment, a graphics-inclusion command, a PDF-page inclusion, a table or a TikZ picture; no image file is copied into this package, the package declares no asset dependency and no image file is redistributed with it. Licence: version arXiv:2502.10353v2 of this article is distributed under the Creative Commons Attribution 4.0 International licence, as recorded in the arXiv licence metadata for this exact version and shown on the abstract page https://arxiv.org/abs/2502.10353v2. A full rights notice is delivered at licenses/arxiv-2502.10353-CC-BY-4.0.txt. Characters: every retained excerpt is pure ASCII, so the delivered engine, XeLaTeX with its default Latin Modern text font, reports no missing character.
\begin{lemma}
    \label{lem:augment}
    Let $\pi$ be a policy that augments the pairwise policy; $\pi(\theta)_{i,j} \geq \pi^{P}(\theta)_{i,j}$ for any $\theta$ for all $i,j$. 
    Let $G = (V,E)$ be a directed graph with $N$ nodes such that nodes $i$ and $i'$ are connected if $\pi(\theta)_{i,v(\theta)_{i'}} = 1$ . 
    Here, $v(\theta)_{i} = j$ if $\pi^{P}(\theta)_{i,j} = 1$.
    Then we have that
    \begin{equation}
        \mathrm{MR}(\pi,\theta,f) = \mathrm{MR}(\pi^{P},\theta,f) 
    \end{equation}
    if each component in $G$ is a complete digraph.  
\end{lemma}

\begin{proof}
% We will first prove the reverse direction; that if $G$ is a complete digraph, then $\pi$ and $\pi^{P}$ have the same match rate. 
Consider patient $i$ in a component that is a connected component of size $k$. 
% This means that $\sum_{j} \pi(\theta)_{i,j} = 1$. 
% We will prove the set of providers unmatched and offered to patient $i$ is non-empty; that is $\lVert \pi(\theta)_{\sigma_{t}} \odot \mathbf{y}^{(t)} \rVert_{1} > 0$ for any ordering $\sigma$ and $\mathbf{y}^{(t)}$ with $\sigma_{t}=i$. 
% $\lVert \pi(\theta)_{\sigma_{t}} \odot \mathbf{y}^{(t)} \rVert_{1} = k$ when $\mathbf{y}^{(t)} = \mathbf{1}$, and decreases by one each time a provider from $\pi(\theta)_{\sigma_{t}}$ is selected. 
For any $j$ with $\pi(\theta)_{i,j}=1$, there exists $k-1$ other $i'$ with $\pi(\theta)_{i',j}=1$ because all patients within the complete graph have provider $j$ on their assortment. 
Therefore, at most $k-1$ other patients can select providers from $\pi(\theta)_{i}$, and therefore $\lVert \pi(\theta)_{\sigma_{t}} \odot \mathbf{y}^{(t)} \rVert_{1}$ can decrease by at most $k-1$ before $\sigma_{t}=i$. 
% Therefore, $\lVert \pi(\theta)_{\sigma_{t}} \odot \mathbf{y}^{(t)} \rVert_{1} \geq k-(k-1) = 1 > 0$. 
Moreover, because $\min(M,N)$ patients have a non-empty assortment, in expectation, we have that $\min(M,N)p$ patients match.%, and so the match rate for the policy $\pi$ is $p \frac{\min(M,N)}{N}$, which is the same as the pairwise policy. 

% Next, we will prove the forward direction: if the match rates are equal between $\pi$ and $\pi^{P}$, then $G$ consists of components that are complete digraphs. 
% Therefore, the only way for $\pi$ to achieve the same match rate as $\pi^{P}$ is to offer $\min(M,N)$ patients a non-empty assortment. 
% Consider some node $u$ in the graph, corresponding to a patient. 
% Suppose that patient $i$ has an assortment of size $k$, indicating that there exist $k$ edges from $i$ to some node. 
% Suppose that there exists a node, $i'$ such that $(i',i)$ is an edge, but $(i,i')$ is not. 
% Then consider the ordering that places $i$ after its neighbors and $i'$. 
% In such an ordering, suppose $i'$ selects $v(\theta)_{i}$ and each neighbor of $i$, $u$, select $v_{u}$. 
% This results in neither $i$ nor its neighbors being available when $i$ must select a patient, leaving an empty assortment.
% Therefore, if patient $i$ has $v_{i'}$ on its assortment, then $i'$ must have $v(\theta)_{i}$ on its assortment for the policy $\pi$. 

% Next, we consider the scenario where there exists a node $w$ such that $w$ is a two-hop neighbor of $i$ but not an immediate neighbor of $i$. 
% Suppose $w$ is adjacent to $i'$, so that $i'$ is on $w$'s assortment. 
% Order the patients such that $w$ comes first, then $i$'s neighbors, then $i$. 
% Let $w$ select $i'$, let $i'$ select $i$, and let all of $i$'s other neighbors select themselves.
% This results in an empty assortment for $i$. 
% Consider the $\min(M,N)$ non-empty assortments; the restriction of $G$ to only these patients has no two-hop neighbors that are not also one-hop neighbors. 
% % In other words, every node in $G$ only has one-hop neighbors.  
% Therefore, in any component, all neighbors are connected, and so $G$ consists of components that are complete digraphs. 
\end{proof}

\end{document}
